\section{Assessment and recommended starting point}

The most promising route beyond \grind{} is not to replace its equality and
arithmetic infrastructure. It is to surround that infrastructure with
proof-structure search and to feed it new, independently validated mathematical
consequences.

Four complementary kinds of progress are available: discovering a certificate,
discovering an invariant, constructing a witness, and --- the extension round's
addition --- computing a \emph{closure} that makes an unbounded family finite. Each of them changes what
is being proved rather than deriving more consequences of a fixed local state,
and that is the distinction the whole architecture rests on. A constructor case
split and an induction principle are not interchangeable --- the former says a
list is empty or a cons, the latter licenses a recursive hypothesis about its
tail. Recognising an existential proposition is not the same as finding a useful
witness for it. Increasing a rewriting budget cannot turn a theorem about one
fixed accumulator into a theorem about all accumulators.

The arithmetic design is deliberately broader than a single SOS tactic, and the
breadth is justified by an actual obstruction rather than by taste. Exact
quadratic decomposition handles hidden global quadratic positivity and is
complete on its fragment. Finite cones exploit useful factors and hypotheses,
which global SOS search structurally cannot do for a domain-dependent statement.
Bernstein trees exploit compact domains. Square-factor and Sturm certificates
handle the univariate zeros that Theorem~\ref{thm:zero} proves subdivision can
never reach. Each of these fails on problems another one decides immediately ---
that is why they belong in one system with a router, and why a failure of one
search is a reason to route rather than a reason to stop.

The closure workers change the shape of that argument rather than adding another
member to the portfolio. A certificate search asks whether a given statement has
a proof in a given language; a closure computes the statement's own consequences
until they stop growing, and then either the target is inside the result or a
counterexample falls out. That is why Section~\ref{sec:closure} can state
decision results and length bounds where Section~\ref{sec:nonlinear} can only
state completeness relative to a dictionary, and it is why a failed closure is
usually a routing signal rather than a request for more budget.

\subsection{What to build first}

\emph{Not} a new SMT solver, and not a rewrite of \grind{}. A narrow Lean
controller that combines automatic generalisation and induction with an exact
witness and certificate pipeline already changes the \emph{shapes} of proofs
that can be found automatically, and it provides clean interfaces for adding
stronger nonlinear and quantified reasoning afterwards. The first useful release
adds induction with generalisation, typed witness construction, and intermediate
certified arithmetic facts. Tightly integrated Boolean search comes after the
scope and explanation contracts are reliable --- which is worth stating plainly,
because the Boolean lane has the most complete prototype here and is still the
right thing to do last.

Two findings should narrow the claims a first release makes. A Lean \code{sos}
project already implements an end-to-end nonlinear-real-arithmetic tactic with
reification, search, rational reconstruction and certificate-based proof
construction; a Lean \code{lp} project already supplies verified linear
programming and a quantified rational-affine fragment. Building either again
would be the second unaudited wrapper, not a contribution
(Remark~\ref{rem:sos-reuse}). The contribution is \emph{when} to request a
certificate, \emph{how} to bind it to the source problem, and \emph{what} to do
with the proved consequence afterwards.

The extension round supplies a smaller and more specific first target than the
design round did, and the recommendation here follows it. The finite-algebra
cover of Section~\ref{sec:covers} has the cheapest checker in this document ---
a bare list of state identifiers, verified by membership, distinctness,
constructor closure and inclusion --- and a counterexample form that is an
ordinary finite tree. It is the right first end-to-end Lean milestone precisely
because it is not the most mathematically interesting lane: the
polynomial-multiplier and telescoping lanes are more substantial and their
source bridges are more delicate. Prove one source theorem through the cheap
lane first, with a recorded axiom inventory, and only then attempt a lane whose
reifier has to get a binomial support or a signed lower index right.

\subsection{Where the risk is}

The largest soundness risk is mistranslation at the Lean/certificate boundary
--- local assumptions, casts, dependent binders, and strict inequalities. The
largest performance risk is generating candidates whose proofs are expensive to
reconstruct. The largest evaluation risk is comparing against an artificially
weak premise set, or leaking target theorems into retrieval.

None of these is mathematical, and that is the useful observation. The hard
parts of this project are not the algorithms; \S\ref{sec:integration} lists them
and they are all engineering. It is worth noticing which way that cuts: a
project whose remaining risk is concentrated in state management and proof
transport is a project whose mathematics is done, and one that should be judged
on whether it ships a narrow correct thing rather than a wide leaky one.

\subsection{What this document has established, and what it has not}

Established: eighteen prototype runs synthesise certificates in fragments
spanning exact polynomial cones, quadratic decomposition, Bernstein subdivision,
univariate Sturm certificates, polynomial and accumulator invariants, affine,
lattice, residue and polynomial witnesses, straight-line integer witness
programs, ground Horn slicing, structural induction with generalisation,
proof-logging CDCL(T), observable-space and target-generated ideal closure,
finite-algebra covers, boundary-safe summation, operator transport, and finite
Kripke countermodels; every stored certificate replays without the search
libraries; the specified mutation and corruption tests are rejected; several
mechanisms solve problems their own restricted precursors cannot; all nine
extension suites re-run in a different environment and reproduce their recorded
counts; and thirteen Lean artefacts elaborate under the pinned v4.34.0
toolchain, seven of them with no axiom dependencies.

Not established: that \forge{} is implemented as a Lean tactic; that the
twenty-one Mathlib-dependent candidate scripts compile; that any certificate
is checked by Lean's kernel; that any checker is correct rather than merely
tested; that the methods outperform any existing Lean tactic; or that any
success rate transfers to user developments. No head-to-head comparison against
\grind{}, Aesop, \code{nlinarith}, LeanHammer, or any other tactic exists
anywhere in this work.

One addition to that second list is specific to the second round and deserves
its own sentence, because it is the claim most likely to be made on this work's
behalf. Elaborating a Lean specimen does not change what the specimen says.
\src{e8}'s indexing sketch now type-checks, and it still proves an equation
between two ordinary-list definitions; it does not close the Church-encoded
query its author declined to claim it closed.

The gap between those two lists is the project. It is not a gap that more
prototype experiments can close, and the second round is the evidence for that:
nine more independent efforts, more certificates, more mutations, more proved
theorems on paper --- and the Lean column still reads \code{NOT\_RUN} in every
one of the eighteen source packages. The gap closes only by writing the Lean
implementation described in Section~\ref{sec:integration} and running the
evaluation described in Section~\ref{sec:evaluation}.

\subsection{The intended end state}

Worth restating, because it constrains the design: the goal is not a permanently
expensive tactic invocation. It is expensive search \emph{once}, then a stable,
readable proof for ordinary recompilation --- a separately proved general lemma
and a short specialisation, with the final theorem not invoking the synthesiser
at all. A proof that succeeds only while a particular search process, solver
build, or model service remains available is an unsatisfactory end product when
an explicit certificate could have been retained instead.

The result is a concrete, falsifiable design for substantially stronger
automation, with working algorithmic building blocks and a precise account of
what still has to be proved and measured. It is more than a list of tactics to
try, and considerably less than a finished Lean prover. Saying which is which,
at every point, is the part of this document that matters most.
